% SPDX-License-Identifier: PMPL-1.0-or-later
% SPDX-FileCopyrightText: 2026 Jonathan D.A. Jewell
%
% The Exceptional Protocols: What We Thought Was Impossible,
% What Actually Was, and an Open Challenge
%
% Intended venue: arXiv preprint (cs.NI / cs.PL / cs.CR)
% Subsequently: HotNets, SIGCOMM Workshop, or USENIX ;login:
%
% Fifth paper in the series:
%   1. Groove Protocol (the mechanism)
%   2. Typing the Boundaries (the boundary taxonomy)
%   3. Graduated Trust (the trust spectrum + IPv6T)
%   4. The Protocol Factory (the bigger vision)
%   5. The Exceptional Protocols (this paper — what can't be done?)

\documentclass[11pt,a4paper]{article}
\usepackage[utf8]{inputenc}
\usepackage[T1]{fontenc}
\usepackage{amsmath,amssymb,amsthm}
\usepackage{listings}
\usepackage{hyperref}
\usepackage{booktabs}
\usepackage{graphicx}
\usepackage{xcolor}
\usepackage[margin=2.5cm]{geometry}

\definecolor{codegreen}{rgb}{0.25,0.49,0.31}
\definecolor{codegray}{rgb}{0.5,0.5,0.5}
\definecolor{codeblue}{rgb}{0.22,0.40,0.72}

\lstset{
  basicstyle=\ttfamily\small,
  keywordstyle=\color{codeblue}\bfseries,
  commentstyle=\color{codegray}\itshape,
  stringstyle=\color{codegreen},
  breaklines=true,
  frame=single,
  captionpos=b
}

\newtheorem{theorem}{Theorem}
\newtheorem{definition}{Definition}
\newtheorem{property}{Property}
\newtheorem{observation}{Observation}

\title{The Exceptional Protocols:\\
What We Thought Was Impossible,\\
What Actually Was, and an Open Challenge}

\author{Jonathan D.A. Jewell\\
\texttt{j.d.a.jewell@open.ac.uk}\\
The Open University}

\date{March 2026}

\begin{document}

\maketitle

\begin{abstract}
During the development of the Groove Protocol and its associated
systems (Typed Frame Router, IPv6T, FLI), we identified six boundary
conditions that appeared to make typed communication impossible. We
present these six ``exceptional protocols'' --- cases where type safety
seemed fundamentally incompatible with the underlying communication
model. We then show how five of the six fell to solutions discovered
during iterative design and testing, leaving a residual set far narrower
than initially expected. We document the two genuinely hard cases that
remain, challenge the community to solve them, and commit to updating
this paper with any progress.
\end{abstract}

\section{Introduction}

Every system has edge cases. When we set out to type every system
boundary --- language, service, protocol, component, runtime, storage
\cite{boundaries2026} --- we expected the edges to be sharp. We expected
large categories of communication to be fundamentally untypeable.

We were wrong about almost all of them.

This paper documents six cases we believed were impossible, the
reasoning that made them seem so, the moment each one broke, and the
solution that emerged. It then documents the two cases that genuinely
remain open, explains why they are harder than the others, and issues
an open challenge to the community.

The purpose is not to claim completeness. It is to be honest about what
we thought, what we learned, and what we still don't know.

\section{Exception 1: The Open Internet}
\label{sec:wan}

\subsection{Why We Thought It Was Impossible}

The Groove Protocol provides compile-time type proofs on localhost: both
endpoints are compiled from the same Idris2 ABI. The proofs hold because
you control both sides.

On the open internet, you don't control the remote endpoint. You didn't
compile their code. The proof chain terminates at your network edge.
Beyond it: strangers, forged headers, untyped bytes. Level 4 (Proven)
seemed restricted to localhost forever.

\subsection{The Moment It Broke}

The assumption was that the internet connects \emph{strangers}. But not
all internet communication is between strangers.

Two users of Burble (a voice platform) are running the same binary,
compiled from the same Idris2 ABI, by the same build pipeline. They are
not strangers to each other's types. They are \emph{proven peers} who
happen to be separated by an ocean.

The internet between them is a dumb pipe. TCP guarantees ordered,
reliable delivery. GRV6 frames carry the type hashes inside the payload,
invisible to all infrastructure. The proofs exist in the binaries at
both ends. The pipe doesn't need to know.

\subsection{The Solution}

The Typed Overlay Network \cite{trust2026}: every panel instance is a
node, every GRV6 connection is an edge. Endpoints compiled from the same
ABI are Level 4 (Proven) to each other regardless of network path.
Different panels sharing the same ABI package for a common capability
are Level 4 for that capability.

\subsection{What Remains}

Communication with true strangers --- endpoints with no shared ABI ---
is capped at Level 3 (Verified) at best, Level 1 (Declared) at minimum.
This is not a limitation of Groove; it is a limitation of reality. You
cannot prove properties of code you haven't seen.

\textbf{Status: Solved for peer-to-peer. Open for stranger-to-stranger.}

\section{Exception 2: Encrypted Tunnels}
\label{sec:encryption}

\subsection{Why We Thought It Was Impossible}

TLS encrypts everything between the handshake and the close. A Typed
Frame Router sitting mid-tunnel cannot inspect the bytes. It cannot
validate type hashes. It cannot filter. The encryption is doing exactly
what it should --- preventing inspection --- and that prevents typing.

We initially listed this as a fundamental limitation: ``types cannot
exist inside encrypted tunnels.''

\subsection{The Moment It Broke}

Types don't need to be \emph{inside} the tunnel. They need to be
\emph{at the endpoints}.

TLS already has a mechanism for this: Application-Layer Protocol
Negotiation (ALPN, RFC~7301). During the TLS handshake --- before any
encrypted application data flows --- both sides declare which protocols
they support. HTTP/2 adoption was built on ALPN.

\subsection{The Solution}

A \texttt{grv6} ALPN token. During the TLS handshake:

\begin{lstlisting}
ClientHello:  ALPN: ["grv6", "h2", "http/1.1"]
ServerHello:  ALPN: "grv6"
\end{lstlisting}

Both sides now know the tunnel carries GRV6 typed frames. After the
handshake, all application data inside the tunnel uses GRV6 framing.
The types were negotiated \emph{before} encryption. No mid-tunnel
inspection needed.

Additionally, TLS encrypts the GRV6 type hashes, preventing traffic
analysis by observers. The encryption that seemed like an obstacle
actually provides a bonus: type privacy.

\subsection{What Remains}

If you are a \emph{middlebox} that needs to inspect typed traffic
passing through an encrypted tunnel you don't terminate, you cannot.
This is by design --- TLS exists to prevent middlebox inspection.
Legitimate cases (corporate DPI, regulatory monitoring) require TLS
termination, at which point the decrypted traffic has GRV6 headers
and typing resumes normally.

\textbf{Status: Solved via ALPN. Mid-tunnel inspection remains
impossible by design.}

\section{Exception 3: Broadcast Protocols}
\label{sec:broadcast}

\subsection{Why We Thought It Was Impossible}

The Typed Frame Router splices TCP connections. Broadcast is
connectionless. There is no stream to splice, no connection to manage,
no point-to-point channel to type. Broadcast seemed to be a
fundamentally different communication model that types couldn't reach.

\subsection{The Moment It Broke}

Groove discovery \emph{is} typed broadcast. It has been from the start.

A system announces its capability manifest via multicast or port
probing. That announcement is a broadcast. The manifest is typed. The
receiver filters by structural type match. If a match is found, the
broadcast upgrades to a connection, and the connection is typed.

We were already doing typed broadcast. We just didn't recognise it as
solving the ``broadcast is impossible'' problem.

\subsection{The Solution}

Two primitives instead of one:

\begin{enumerate}
\item \textbf{Typed Broadcast Filter}: embeds type hashes in broadcast
  payloads. Works on mDNS (TXT records, 255 bytes), BLE advertising
  (manufacturer data, 31 bytes --- almost exactly one SHA-256 hash),
  WiFi beacons (vendor-specific IEs, 255 bytes), and Zigbee
  (manufacturer-specific clusters).

\item \textbf{Typed Frame Router}: once the broadcast discovers a match,
  the connection is established and the TFR handles the typed stream.
\end{enumerate}

The lifecycle: broadcast $\rightarrow$ discover $\rightarrow$
connect $\rightarrow$ communicate. Broadcast is the typed discovery
layer. Connection is the typed data layer.

\subsection{What Remains}

Pure broadcast with no connection upgrade --- systems that broadcast
data continuously without any listener establishing a connection (e.g.,
a lighthouse broadcasting a GPS-timestamped beacon). These can carry
type hashes in the broadcast payload, but the receiver has no way to
confirm type conformance. The hash is an assertion in the wind.

\textbf{Status: Solved for discovery-then-connect. Open for
fire-and-forget broadcast.}

\section{Exception 4: IPv4-to-IPv6 Translation}
\label{sec:ipv4v6}

\subsection{Why We Thought It Was Limited}

Our initial IPv4$\rightarrow$IPv6 transition design used HTTP-level
upgrade headers: accept an IPv4 connection, respond with
\texttt{Groove-Upgrade-To: ipv6}, let the client reconnect on IPv6.
Two round trips, state loss, client awareness required.

This seemed like the only option: the application layer is where Groove
operates, so the upgrade has to happen at the application layer.

\subsection{The Moment It Broke}

``Why not strip it at the frame level instead of the packet level?''

TCP streams have frames (segments). The Zig FFI can accept a TCP
connection on IPv4 and open a corresponding connection on IPv6, then
\texttt{splice(2)} the bytes bidirectionally. Zero-copy. Transparent.
The client never knows. The server only sees IPv6.

\subsection{The Solution}

The Typed Frame Router operating as an IPv4$\rightarrow$IPv6 proxy at
TCP Layer 4. On Linux, \texttt{splice(2)} transfers bytes between socket
file descriptors without copying to userspace. The proxy is approximately
60 lines of Zig. It is formally verified in Idris2 (transport
transparency, bounded memory, direction safety, liveness).

It is now the 9th core primitive in the proven-servers repository.

\subsection{What Remains}

Nothing. This is completely solved.

\textbf{Status: Solved. Formally verified. Tests passing.}

\section{Exception 5: Browser-to-Native Boundary}
\label{sec:browser}

\subsection{Why We Thought It Was Fragile}

Browser extensions talk to local services via \texttt{fetch()} to
localhost. The request is untyped HTTP. The response is parsed JSON.
Neither side has compile-time proof of the other's types. The browser
security sandbox actively prevents the kind of low-level socket access
that the Typed Frame Router uses.

\subsection{The Moment It Broke}

The Groove harness pattern: a browser extension that probes localhost
for Groove manifests using standard \texttt{fetch()} and establishes
type-safe connections. The types cross the boundary via the chain:
Idris2 ABI $\rightarrow$ C headers $\rightarrow$ Zig FFI $\rightarrow$
HTTP $\rightarrow$ \texttt{fetch()} $\rightarrow$ JavaScript.

The browser sandbox doesn't prevent \texttt{fetch()} to localhost.
1Password, VS Code, Figma, and Obsidian all use this pattern. The
precedent is established.

\subsection{The Solution}

The Groove Browser Harness: a browser extension running the same
\texttt{GrooveHarness} constructor as a Gossamer panel, but inside a
service worker. CORS headers on groove endpoints enable the
\texttt{fetch()}. Dual-stack probing (\texttt{[::1]} first,
\texttt{127.0.0.1} fallback) handles all configurations.

Firefox embraces it (persistent background, open web philosophy).
Chrome accepts it with Manifest V3 friction (service worker lifecycle).
Safari tries to restrict it (App Transport Security, App Store
gatekeeping).

\subsection{What Remains}

Safari. Apple's incentive is to prevent exactly this kind of
disintermediation. Whether this is a technical limitation or a business
decision is a question for Apple.

\textbf{Status: Solved for Firefox and Chrome. Obstructed (not
impossible) on Safari.}

\section{Exception 6: Hardware-Accelerated Bypass}
\label{sec:hwbypass}

\subsection{Why We Thought It Was Impossible}

RDMA (Remote Direct Memory Access), DPDK (Data Plane Development Kit),
and SR-IOV (Single Root I/O Virtualisation) bypass the kernel entirely.
Network packets go directly from the NIC to userspace memory (or to a
virtual machine's memory) without the kernel ever seeing them.

The Typed Frame Router uses \texttt{splice(2)}, which is a kernel
operation. If the kernel never sees the packets, \texttt{splice(2)}
cannot intercept them. The TFR is blind.

\subsection{Why It's Still Hard}

This is genuinely difficult because the bypass is the \emph{point}.
RDMA exists because kernel processing is too slow for high-frequency
trading, HPC, and storage networks. Any interception --- even a single
branch to check a type hash --- adds latency that defeats the purpose.

\subsection{Possible Approaches (Speculative)}

\begin{enumerate}
\item \textbf{Userspace typed filter}: DPDK provides a userspace network
  stack. A typed frame filter could run inside the DPDK pipeline,
  checking type hashes as part of the existing packet processing loop.
  This adds computation but no kernel round-trip. The overhead would be
  a single SHA-256 comparison (one cache line) per packet.

\item \textbf{SmartNIC offload}: modern SmartNICs (NVIDIA BlueField,
  AMD Pensando) can run custom programs on the NIC itself. A typed frame
  filter running on the SmartNIC would intercept packets \emph{before}
  they reach the host, with zero host CPU overhead. The NIC becomes a
  typed boundary.

\item \textbf{eBPF at the XDP layer}: Linux Express Data Path (XDP)
  runs eBPF programs on the NIC driver, before the kernel network stack.
  A GRV6 magic check in XDP ($< 10$ ns) could route typed frames to a
  validation path while passing untyped frames through at full speed.
\end{enumerate}

None of these are implemented. Each requires specialised hardware or
kernel access. We list them as directions, not solutions.

\textbf{Status: Open. Genuinely hard. Challenge issued.}

\section{Exception 7: Unknown Strangers}
\label{sec:strangers}

\subsection{Why It's Impossible}

A packet arrives from an IP address you've never seen. It has no GRV6
header. No Groove manifest. No type hash. No ALPN token. It is raw
bytes from an unknown source.

You cannot type this. You cannot validate what was never declared. You
cannot prove properties of code you've never seen. This is not a
limitation of Groove or GRV6 or any typed system. It is a limitation of
information: you cannot reason about what you don't have.

\subsection{Why It Will Always Be Impossible}

Typing requires at minimum a \emph{declaration} from the sender: ``I
claim this data is of type X.'' Without that declaration, there is
nothing to check, nothing to match, nothing to prove. The sender must
participate.

This is analogous to encryption: you cannot encrypt a conversation
if the other party refuses to exchange keys. You cannot type a
conversation if the other party refuses to declare types.

\subsection{What You CAN Do}

\begin{enumerate}
\item \textbf{Reject}: refuse untyped traffic entirely. This is
  appropriate for systems that only communicate with known peers (the
  typed overlay).
\item \textbf{Quarantine}: accept untyped traffic but process it in a
  restricted sandbox. Treat it as Level 0 (Untyped) and apply maximum
  runtime validation.
\item \textbf{Upgrade}: if the untyped traffic follows a known protocol
  (HTTP, gRPC), wrap it in a GRV6 frame at the boundary. A Groove
  Bridge does exactly this. The bridge is the translation layer between
  the untyped internet and the typed overlay.
\end{enumerate}

\textbf{Status: Impossible. Fundamentally. Not a challenge --- a
physical law of information.}

\section{The Scorecard}

\begin{table}[h]
\centering
\small
\begin{tabular}{clcl}
\toprule
\# & Exception & Status & Solution \\
\midrule
1 & Open Internet (WAN) & \textbf{Solved} & Typed overlay (shared ABI = Level 4) \\
2 & Encrypted Tunnels & \textbf{Solved} & TLS ALPN negotiation \\
3 & Broadcast Protocols & \textbf{Solved} & Typed Broadcast Filter + upgrade \\
4 & IPv4$\rightarrow$IPv6 & \textbf{Solved} & Frame-level splice (verified) \\
5 & Browser$\leftrightarrow$Native & \textbf{Solved} & Groove Browser Harness \\
6 & Hardware Bypass & \textbf{Open} & DPDK/SmartNIC/XDP (speculative) \\
7 & Unknown Strangers & \textbf{Impossible} & Information-theoretic limit \\
\bottomrule
\end{tabular}
\caption{Exceptional protocol scorecard}
\label{tab:scorecard}
\end{table}

Of seven exceptions we identified, five were solved during iterative
design and testing. One remains open but has speculative approaches. One
is fundamentally impossible.

The genuinely untypeable surface is far narrower than we expected:
hardware-accelerated kernel bypass (a niche use case) and packets from
strangers who refuse to participate (an information-theoretic limit, not
a technical one).

\section{An Open Challenge}

We challenge the community on Exception 6 (Hardware-Accelerated Bypass):

\begin{quote}
\emph{Can a GRV6 type hash be validated at line rate on a hardware-
accelerated data path (RDMA, DPDK, or SR-IOV) without defeating the
purpose of the acceleration? Specifically: can the overhead be bounded
to $< 100$ ns per packet, which is within the jitter budget of most
HPC and HFT applications?}
\end{quote}

We believe the answer is yes, via SmartNIC offload or XDP/eBPF, but we
have not demonstrated it. We commit to updating this paper --- with a
new section, not a revision --- when we or anyone else produces a
working implementation.

Exception 7 (Unknown Strangers) is not a challenge. It is a boundary.
You cannot type what was never declared. This is as fundamental as ``you
cannot decrypt what was never encrypted.'' We state it not as a
limitation to be overcome but as a condition of reality to be understood.

\section{What We Learned}

The most important lesson from this exercise is that the boundary between
``possible'' and ``impossible'' was in a different place than we thought.

We began by listing six things that seemed impossible. We ended with two
--- and one of those is not ``impossible'' but ``hard.'' The apparently
impossible cases fell not because we found clever workarounds but because
we had framed the problems wrong:

\begin{itemize}
\item We said ``you can't type the internet'' when we meant ``you can't
  type strangers.'' The internet between peers is just a wire.
\item We said ``you can't type encrypted tunnels'' when we meant ``you
  can't type mid-tunnel.'' The endpoints are unencrypted.
\item We said ``you can't type broadcast'' when we meant ``you can't
  splice broadcast.'' You can filter it.
\item We said ``you can't type the browser boundary'' when we meant ``the
  browser sandbox prevents socket access.'' It doesn't prevent
  \texttt{fetch()}.
\end{itemize}

In each case, the impossibility was an artefact of how we defined the
problem, not a property of the problem itself. When we reframed the
question --- from ``can we type this boundary?'' to ``where exactly in
this boundary do the types die, and can we move that point?'' --- the
answers appeared.

We do not claim to have typed the world. We claim to have found that
the world is more typeable than it looks, and that the untypeable
residue is smaller, and more precisely bounded, than we expected.

\section{Conclusion}

Five out of seven. That is the score.

Five communication patterns that we believed were fundamentally
incompatible with type safety turned out to be solvable --- not through
heroic engineering but through precise reframing of where the types
need to exist.

One remains open: hardware-accelerated kernel bypass. We challenge the
community to solve it. We believe it can be done at $< 100$ ns per
packet, and we will report any progress.

One is impossible: typing traffic from strangers who refuse to
participate. We do not challenge anyone to solve this. We document it
as a boundary condition of typed communication, as fundamental as the
impossibility of encrypting a conversation with an unwilling
participant.

The world is more typeable than it looks. The exceptions are
exceptional. And the challenge is open.

\bigskip
\noindent
\textbf{Update commitment:} This paper will be updated --- not revised,
but extended with new sections --- when progress is made on Exception 6.
Updates will be clearly dated and attributed. The original text will not
be modified.

\bibliographystyle{plain}
\begin{thebibliography}{10}

\bibitem{groove2026}
J.~D.~A.~Jewell.
\newblock Groove: Formally Verified System Composition via Dependent Types and
  Linear Resources.
\newblock arXiv preprint, 2026.

\bibitem{boundaries2026}
J.~D.~A.~Jewell.
\newblock Typing the Boundaries: Preserving Type Safety Across Every System
  Crossing Point.
\newblock arXiv preprint, 2026.

\bibitem{trust2026}
J.~D.~A.~Jewell.
\newblock Graduated Trust: A Spectrum of Type Safety from Localhost to the Open
  Internet.
\newblock arXiv preprint, 2026.

\bibitem{factory2026}
J.~D.~A.~Jewell.
\newblock The Protocol Factory: Malleable, Ductile Protocol Creation via
  Formally Verified Typed Overlays.
\newblock arXiv preprint, 2026.

\bibitem{rfc7301}
S.~Friedl, A.~Popov, A.~Langley, and E.~Stephan.
\newblock Transport Layer Security (TLS) Application-Layer Protocol Negotiation
  Extension.
\newblock RFC 7301, 2014.

\end{thebibliography}

\end{document}
